% Energy-price fork (bow-tie): one electricity price reaches the economy two ways -- producers' costs
% rise (a broad approximation) and households pay more directly (represented natively), both converging
% on households. Plain language, trust tags (size reliable vs direction only), no model symbols/code.
\documentclass[border=16pt]{standalone}
% Shared preamble for every standalone diagram. Each diagram does:
%   \documentclass[border=16pt]{standalone}
%   % Shared preamble for every standalone diagram. Each diagram does:
%   \documentclass[border=16pt]{standalone}
%   % Shared preamble for every standalone diagram. Each diagram does:
%   \documentclass[border=16pt]{standalone}
%   \input{_preamble}
%   \begin{document}\begin{tikzpicture} ... \end{tikzpicture}\end{document}
% Compiled once to a tight-cropped PDF (the reusable asset); the Beamer deck then
% \includegraphics those PDFs. Keeps diagram source out of the deck preamble.
%
% Visual language: FT / Nature / Quanta editorial discipline -- a warm paper ground,
% a serif declarative title over a grey dek, restrained colour where HUE = identity
% (sign carried by glyph + line style, so it survives greyscale and colour blindness).

\usepackage[T1]{fontenc}
\usepackage[scaled=1.0]{helvet}          % editorial sans (matches style.py fallback stack)
\usepackage[scaled=1.04]{XCharter}       % transitional serif -- declarative TITLES only (FT/Quanta signature)
\renewcommand{\familydefault}{\sfdefault}% sans stays the default; the serif is reached via \rmfamily
\usepackage{microtype}                   % provides \textls for the letter-spaced kicker
\usepackage{amsmath}
\usepackage{amssymb}
\usepackage{fontawesome5}
\usepackage{tikz}
\usetikzlibrary{positioning, arrows.meta, calc, fit, backgrounds,
                shapes.geometric, shapes.arrows, shapes.misc,
                decorations.pathreplacing, decorations.pathmorphing}

\input{../theme/ogclews-colors}

% ---- editorial canvas (FT/Quanta): warm paper ground + a context grey for inactive structure ----
\definecolor{paper}{HTML}{FBF7F0}        % warm canvas, lighter than FT cream so navy/teal keep full contrast
\definecolor{context}{HTML}{CBC3B8}      % warm grey -- de-emphasised / inactive structure ONLY (never an identity)
\providecommand{\ogground}{paper}        % a Nature/journal cut flips to white: \def\ogground{white} before \input
\pagecolor{\ogground}
\tikzset{every picture/.style={line cap=round, line join=round}}  % soft, modern strokes everywhere

% Narrow cards: suppress mid-word hyphenation; let lines break at spaces (ragged), with a little
% stretch so short card lines do not overflow. Keeps the figures clean (no "sys-tem" breaks).
\hyphenpenalty=10000 \exhyphenpenalty=10000 \emergencystretch=2em \hbadness=10000

% ---- shared semantic styles -------------------------------------------------
\tikzset{
  % a model / container panel: faint identity tint, colored border, bold colored title set inline
  panel/.style={rounded corners=4pt, line width=1.0pt, inner sep=13pt,
                align=left, fill=white},
  ogpanel/.style={panel, draw=ogc, fill=ogc!6},
  clewspanel/.style={panel, draw=clews, fill=clews!6},
  policypanel/.style={panel, draw=policy, fill=policy!6},
  card/.style={rounded corners=3pt, line width=0.8pt, inner sep=10pt, align=left,
               draw=grid, fill=white},
  % de-emphasised / inactive structure (grey-for-context, never a recoloured identity)
  ghostcard/.style={card, draw=context, dashed, fill=paper},
  % small uppercase pill used for direction / theory-status tags
  pill/.style={rounded corners=2.5pt, inner xsep=5pt, inner ysep=2.5pt,
               font=\scriptsize\bfseries, text=white},
  % flows -- slimmer, cleaner arrowheads (FT restraint)
  flow/.style={-{Stealth[length=2.6mm,width=2.0mm]}, line width=1.4pt},
  thinflow/.style={-{Stealth[length=2.2mm,width=1.7mm]}, line width=1.0pt},
  ghostflow/.style={thinflow, draw=context, dashed},     % scaffold / inactive direction
  seamflow/.style={flow, dashed},                        % stubbed / illustrative seam
  % annotation callout
  note/.style={font=\footnotesize\itshape, text=sub, align=center},
  kicklabel/.style={font=\scriptsize\bfseries, text=claret},
}

% kicker rule + label, used to title a standalone diagram lightly (optional)
\newcommand{\kicker}[2]{% #1 = color, #2 = TEXT
  \tikz[baseline]{\draw[line width=2pt,#1] (0,0.18) -- (0.55,0.18);}\,%
  {\scriptsize\bfseries\color{#1} #2}}

% reusable editorial header (FT/Quanta): claret tick + letter-spaced kicker + serif title + grey dek,
% all left-aligned to a shared margin.  \fighead{x}{y (top-left)}{KICKER}{Declarative title}{dek}
\newcommand{\fighead}[5]{%
  \draw[claret, line width=2pt] (#1,#2) -- ++(0.62,0);
  \node[anchor=north west, inner sep=0pt, font=\scriptsize\bfseries, text=claret]
       (fhk) at ($(#1,#2)+(0,-0.18)$) {\textls[140]{#3}};
  \node[anchor=north west, inner sep=0pt, font=\rmfamily\LARGE\bfseries, text=ink]
       (fht) at ($(fhk.south west)+(0,-0.13)$) {#4};
  \node[anchor=north west, inner sep=0pt, font=\small, text=sub]
       (fhd) at ($(fht.south west)+(0,-0.13)$) {#5};
}
% bottom-left source / credit line (the FT/Nature/Quanta honesty finish). \figsource{x}{y}{text}
\newcommand{\figsource}[3]{%
  \node[anchor=north west, inner sep=0pt, font=\scriptsize, text=mute, align=left] at (#1,#2) {#3};
}
% a filled tag set as a card's first line. \tagline{fill colour}{TEXT}  (shared across figures)
\newcommand{\tagline}[2]{{\colorbox{#1}{\textcolor{white}{\scriptsize\bfseries #2}}}\\[4pt]}

% ---- plain-language causal-chain flow (general audience, icon-led) -----------
\tikzset{
  stagecircle/.style={circle, minimum size=1.55cm, inner sep=0pt, text=white,
                      font=\LARGE, line width=0pt},
  stagelabel/.style={font=\small, text=ink, align=center, text width=3.1cm},
  flowedge/.style={-{Stealth[length=3.4mm,width=3mm]}, line width=1.7pt, sub},
  flowtitle/.style={font=\Large\bfseries, text=ink},
  flowcap/.style={font=\small, text=sub, align=center},
}
% \stage{id}{x}{y}{fillcolor}{icon}{label}
\newcommand{\stage}[6]{%
  \node[stagecircle, fill=#4] (#1) at (#2,#3) {#5};
  \node[stagelabel, below=2.5mm of #1] {#6};
}

% a small "data" sparkline card that shows the by-X profile a channel uses, with a
% source tag -- the visible link between a mechanism and its model effect.
% \datacard{x}{y}{title}{x-axis label}{source tag}{colour}{plot coords}
% coords are relative to (x,y-0.05), spanning x in [-1.1,1.1], y in [-0.4,0.45].
\newcommand{\datacard}[7]{%
  \node[draw=grid, rounded corners=3pt, fill=white, minimum width=2.8cm,
        minimum height=1.8cm, inner sep=0pt] at (#1,#2) {};
  \node[font=\tiny\bfseries, text=sub, anchor=north] at (#1,#2+0.80) {#3};
  \begin{scope}[shift={(#1,#2-0.05)}]
    \draw[grid, line width=0.6pt] (-1.12,-0.42) -- (1.12,-0.42);
    \draw[#6, line width=1.4pt, line cap=round] plot[smooth] coordinates {#7};
    \node[font=\tiny, text=mute, anchor=north west] at (-1.08,-0.42) {#4};
  \end{scope}
  \node[font=\tiny\bfseries, text=white, fill=#6, rounded corners=1.5pt, inner xsep=2.5pt,
        inner ysep=0.8pt, anchor=south east] at (#1+1.3,#2-0.86) {#5};
}

%   \begin{document}\begin{tikzpicture} ... \end{tikzpicture}\end{document}
% Compiled once to a tight-cropped PDF (the reusable asset); the Beamer deck then
% \includegraphics those PDFs. Keeps diagram source out of the deck preamble.
%
% Visual language: FT / Nature / Quanta editorial discipline -- a warm paper ground,
% a serif declarative title over a grey dek, restrained colour where HUE = identity
% (sign carried by glyph + line style, so it survives greyscale and colour blindness).

\usepackage[T1]{fontenc}
\usepackage[scaled=1.0]{helvet}          % editorial sans (matches style.py fallback stack)
\usepackage[scaled=1.04]{XCharter}       % transitional serif -- declarative TITLES only (FT/Quanta signature)
\renewcommand{\familydefault}{\sfdefault}% sans stays the default; the serif is reached via \rmfamily
\usepackage{microtype}                   % provides \textls for the letter-spaced kicker
\usepackage{amsmath}
\usepackage{amssymb}
\usepackage{fontawesome5}
\usepackage{tikz}
\usetikzlibrary{positioning, arrows.meta, calc, fit, backgrounds,
                shapes.geometric, shapes.arrows, shapes.misc,
                decorations.pathreplacing, decorations.pathmorphing}

% OG-Core x CLEWS — shared color palette.
% Single source of truth, mirrored from ogclews_link/style.py (the matplotlib figure
% theme) so the TikZ diagrams, the Beamer deck, and the data figures all share one
% editorial, colorblind-safe palette. \input this from any preamble; it only defines
% colors (no package loads), so it is safe in both standalone and beamer documents.

\definecolor{ink}{HTML}{222222}      % primary text / heavy rules
\definecolor{sub}{HTML}{555555}      % secondary text
\definecolor{mute}{HTML}{888888}     % tertiary / source lines
\definecolor{grid}{HTML}{E6E6E6}     % faint rules, fills

% the two models + policy — the load-bearing semantic colors
\definecolor{ogc}{HTML}{0F5499}      % OG-Core      (FT oxford blue)
\definecolor{clews}{HTML}{0D7680}    % CLEWS/OSeMOSYS (teal)
\definecolor{policy}{HTML}{E3120B}   % policy lever (Economist red)

% accents / signed values (match style.py)
\definecolor{claret}{HTML}{990F3D}   % kicker rule
\definecolor{gain}{HTML}{2166AC}     % positive change (blue)
\definecolor{loss}{HTML}{B2182B}     % negative change (red)
\definecolor{amber}{HTML}{E69F00}    % categorical 4
\definecolor{purple}{HTML}{7A6FAC}   % categorical 5
\definecolor{slate}{HTML}{7F8C8D}    % categorical 6 / neutral

% sequential blue ramp (ordered income groups, poorest -> richest)
\definecolor{seqA}{HTML}{C6DBEF}
\definecolor{seqB}{HTML}{9ECAE1}
\definecolor{seqC}{HTML}{6BAED6}
\definecolor{seqD}{HTML}{4292C6}
\definecolor{seqE}{HTML}{2171B5}
\definecolor{seqF}{HTML}{08519C}
\definecolor{seqG}{HTML}{08306B}


% ---- editorial canvas (FT/Quanta): warm paper ground + a context grey for inactive structure ----
\definecolor{paper}{HTML}{FBF7F0}        % warm canvas, lighter than FT cream so navy/teal keep full contrast
\definecolor{context}{HTML}{CBC3B8}      % warm grey -- de-emphasised / inactive structure ONLY (never an identity)
\providecommand{\ogground}{paper}        % a Nature/journal cut flips to white: \def\ogground{white} before \input
\pagecolor{\ogground}
\tikzset{every picture/.style={line cap=round, line join=round}}  % soft, modern strokes everywhere

% Narrow cards: suppress mid-word hyphenation; let lines break at spaces (ragged), with a little
% stretch so short card lines do not overflow. Keeps the figures clean (no "sys-tem" breaks).
\hyphenpenalty=10000 \exhyphenpenalty=10000 \emergencystretch=2em \hbadness=10000

% ---- shared semantic styles -------------------------------------------------
\tikzset{
  % a model / container panel: faint identity tint, colored border, bold colored title set inline
  panel/.style={rounded corners=4pt, line width=1.0pt, inner sep=13pt,
                align=left, fill=white},
  ogpanel/.style={panel, draw=ogc, fill=ogc!6},
  clewspanel/.style={panel, draw=clews, fill=clews!6},
  policypanel/.style={panel, draw=policy, fill=policy!6},
  card/.style={rounded corners=3pt, line width=0.8pt, inner sep=10pt, align=left,
               draw=grid, fill=white},
  % de-emphasised / inactive structure (grey-for-context, never a recoloured identity)
  ghostcard/.style={card, draw=context, dashed, fill=paper},
  % small uppercase pill used for direction / theory-status tags
  pill/.style={rounded corners=2.5pt, inner xsep=5pt, inner ysep=2.5pt,
               font=\scriptsize\bfseries, text=white},
  % flows -- slimmer, cleaner arrowheads (FT restraint)
  flow/.style={-{Stealth[length=2.6mm,width=2.0mm]}, line width=1.4pt},
  thinflow/.style={-{Stealth[length=2.2mm,width=1.7mm]}, line width=1.0pt},
  ghostflow/.style={thinflow, draw=context, dashed},     % scaffold / inactive direction
  seamflow/.style={flow, dashed},                        % stubbed / illustrative seam
  % annotation callout
  note/.style={font=\footnotesize\itshape, text=sub, align=center},
  kicklabel/.style={font=\scriptsize\bfseries, text=claret},
}

% kicker rule + label, used to title a standalone diagram lightly (optional)
\newcommand{\kicker}[2]{% #1 = color, #2 = TEXT
  \tikz[baseline]{\draw[line width=2pt,#1] (0,0.18) -- (0.55,0.18);}\,%
  {\scriptsize\bfseries\color{#1} #2}}

% reusable editorial header (FT/Quanta): claret tick + letter-spaced kicker + serif title + grey dek,
% all left-aligned to a shared margin.  \fighead{x}{y (top-left)}{KICKER}{Declarative title}{dek}
\newcommand{\fighead}[5]{%
  \draw[claret, line width=2pt] (#1,#2) -- ++(0.62,0);
  \node[anchor=north west, inner sep=0pt, font=\scriptsize\bfseries, text=claret]
       (fhk) at ($(#1,#2)+(0,-0.18)$) {\textls[140]{#3}};
  \node[anchor=north west, inner sep=0pt, font=\rmfamily\LARGE\bfseries, text=ink]
       (fht) at ($(fhk.south west)+(0,-0.13)$) {#4};
  \node[anchor=north west, inner sep=0pt, font=\small, text=sub]
       (fhd) at ($(fht.south west)+(0,-0.13)$) {#5};
}
% bottom-left source / credit line (the FT/Nature/Quanta honesty finish). \figsource{x}{y}{text}
\newcommand{\figsource}[3]{%
  \node[anchor=north west, inner sep=0pt, font=\scriptsize, text=mute, align=left] at (#1,#2) {#3};
}
% a filled tag set as a card's first line. \tagline{fill colour}{TEXT}  (shared across figures)
\newcommand{\tagline}[2]{{\colorbox{#1}{\textcolor{white}{\scriptsize\bfseries #2}}}\\[4pt]}

% ---- plain-language causal-chain flow (general audience, icon-led) -----------
\tikzset{
  stagecircle/.style={circle, minimum size=1.55cm, inner sep=0pt, text=white,
                      font=\LARGE, line width=0pt},
  stagelabel/.style={font=\small, text=ink, align=center, text width=3.1cm},
  flowedge/.style={-{Stealth[length=3.4mm,width=3mm]}, line width=1.7pt, sub},
  flowtitle/.style={font=\Large\bfseries, text=ink},
  flowcap/.style={font=\small, text=sub, align=center},
}
% \stage{id}{x}{y}{fillcolor}{icon}{label}
\newcommand{\stage}[6]{%
  \node[stagecircle, fill=#4] (#1) at (#2,#3) {#5};
  \node[stagelabel, below=2.5mm of #1] {#6};
}

% a small "data" sparkline card that shows the by-X profile a channel uses, with a
% source tag -- the visible link between a mechanism and its model effect.
% \datacard{x}{y}{title}{x-axis label}{source tag}{colour}{plot coords}
% coords are relative to (x,y-0.05), spanning x in [-1.1,1.1], y in [-0.4,0.45].
\newcommand{\datacard}[7]{%
  \node[draw=grid, rounded corners=3pt, fill=white, minimum width=2.8cm,
        minimum height=1.8cm, inner sep=0pt] at (#1,#2) {};
  \node[font=\tiny\bfseries, text=sub, anchor=north] at (#1,#2+0.80) {#3};
  \begin{scope}[shift={(#1,#2-0.05)}]
    \draw[grid, line width=0.6pt] (-1.12,-0.42) -- (1.12,-0.42);
    \draw[#6, line width=1.4pt, line cap=round] plot[smooth] coordinates {#7};
    \node[font=\tiny, text=mute, anchor=north west] at (-1.08,-0.42) {#4};
  \end{scope}
  \node[font=\tiny\bfseries, text=white, fill=#6, rounded corners=1.5pt, inner xsep=2.5pt,
        inner ysep=0.8pt, anchor=south east] at (#1+1.3,#2-0.86) {#5};
}

%   \begin{document}\begin{tikzpicture} ... \end{tikzpicture}\end{document}
% Compiled once to a tight-cropped PDF (the reusable asset); the Beamer deck then
% \includegraphics those PDFs. Keeps diagram source out of the deck preamble.
%
% Visual language: FT / Nature / Quanta editorial discipline -- a warm paper ground,
% a serif declarative title over a grey dek, restrained colour where HUE = identity
% (sign carried by glyph + line style, so it survives greyscale and colour blindness).

\usepackage[T1]{fontenc}
\usepackage[scaled=1.0]{helvet}          % editorial sans (matches style.py fallback stack)
\usepackage[scaled=1.04]{XCharter}       % transitional serif -- declarative TITLES only (FT/Quanta signature)
\renewcommand{\familydefault}{\sfdefault}% sans stays the default; the serif is reached via \rmfamily
\usepackage{microtype}                   % provides \textls for the letter-spaced kicker
\usepackage{amsmath}
\usepackage{amssymb}
\usepackage{fontawesome5}
\usepackage{tikz}
\usetikzlibrary{positioning, arrows.meta, calc, fit, backgrounds,
                shapes.geometric, shapes.arrows, shapes.misc,
                decorations.pathreplacing, decorations.pathmorphing}

% OG-Core x CLEWS — shared color palette.
% Single source of truth, mirrored from ogclews_link/style.py (the matplotlib figure
% theme) so the TikZ diagrams, the Beamer deck, and the data figures all share one
% editorial, colorblind-safe palette. \input this from any preamble; it only defines
% colors (no package loads), so it is safe in both standalone and beamer documents.

\definecolor{ink}{HTML}{222222}      % primary text / heavy rules
\definecolor{sub}{HTML}{555555}      % secondary text
\definecolor{mute}{HTML}{888888}     % tertiary / source lines
\definecolor{grid}{HTML}{E6E6E6}     % faint rules, fills

% the two models + policy — the load-bearing semantic colors
\definecolor{ogc}{HTML}{0F5499}      % OG-Core      (FT oxford blue)
\definecolor{clews}{HTML}{0D7680}    % CLEWS/OSeMOSYS (teal)
\definecolor{policy}{HTML}{E3120B}   % policy lever (Economist red)

% accents / signed values (match style.py)
\definecolor{claret}{HTML}{990F3D}   % kicker rule
\definecolor{gain}{HTML}{2166AC}     % positive change (blue)
\definecolor{loss}{HTML}{B2182B}     % negative change (red)
\definecolor{amber}{HTML}{E69F00}    % categorical 4
\definecolor{purple}{HTML}{7A6FAC}   % categorical 5
\definecolor{slate}{HTML}{7F8C8D}    % categorical 6 / neutral

% sequential blue ramp (ordered income groups, poorest -> richest)
\definecolor{seqA}{HTML}{C6DBEF}
\definecolor{seqB}{HTML}{9ECAE1}
\definecolor{seqC}{HTML}{6BAED6}
\definecolor{seqD}{HTML}{4292C6}
\definecolor{seqE}{HTML}{2171B5}
\definecolor{seqF}{HTML}{08519C}
\definecolor{seqG}{HTML}{08306B}


% ---- editorial canvas (FT/Quanta): warm paper ground + a context grey for inactive structure ----
\definecolor{paper}{HTML}{FBF7F0}        % warm canvas, lighter than FT cream so navy/teal keep full contrast
\definecolor{context}{HTML}{CBC3B8}      % warm grey -- de-emphasised / inactive structure ONLY (never an identity)
\providecommand{\ogground}{paper}        % a Nature/journal cut flips to white: \def\ogground{white} before \input
\pagecolor{\ogground}
\tikzset{every picture/.style={line cap=round, line join=round}}  % soft, modern strokes everywhere

% Narrow cards: suppress mid-word hyphenation; let lines break at spaces (ragged), with a little
% stretch so short card lines do not overflow. Keeps the figures clean (no "sys-tem" breaks).
\hyphenpenalty=10000 \exhyphenpenalty=10000 \emergencystretch=2em \hbadness=10000

% ---- shared semantic styles -------------------------------------------------
\tikzset{
  % a model / container panel: faint identity tint, colored border, bold colored title set inline
  panel/.style={rounded corners=4pt, line width=1.0pt, inner sep=13pt,
                align=left, fill=white},
  ogpanel/.style={panel, draw=ogc, fill=ogc!6},
  clewspanel/.style={panel, draw=clews, fill=clews!6},
  policypanel/.style={panel, draw=policy, fill=policy!6},
  card/.style={rounded corners=3pt, line width=0.8pt, inner sep=10pt, align=left,
               draw=grid, fill=white},
  % de-emphasised / inactive structure (grey-for-context, never a recoloured identity)
  ghostcard/.style={card, draw=context, dashed, fill=paper},
  % small uppercase pill used for direction / theory-status tags
  pill/.style={rounded corners=2.5pt, inner xsep=5pt, inner ysep=2.5pt,
               font=\scriptsize\bfseries, text=white},
  % flows -- slimmer, cleaner arrowheads (FT restraint)
  flow/.style={-{Stealth[length=2.6mm,width=2.0mm]}, line width=1.4pt},
  thinflow/.style={-{Stealth[length=2.2mm,width=1.7mm]}, line width=1.0pt},
  ghostflow/.style={thinflow, draw=context, dashed},     % scaffold / inactive direction
  seamflow/.style={flow, dashed},                        % stubbed / illustrative seam
  % annotation callout
  note/.style={font=\footnotesize\itshape, text=sub, align=center},
  kicklabel/.style={font=\scriptsize\bfseries, text=claret},
}

% kicker rule + label, used to title a standalone diagram lightly (optional)
\newcommand{\kicker}[2]{% #1 = color, #2 = TEXT
  \tikz[baseline]{\draw[line width=2pt,#1] (0,0.18) -- (0.55,0.18);}\,%
  {\scriptsize\bfseries\color{#1} #2}}

% reusable editorial header (FT/Quanta): claret tick + letter-spaced kicker + serif title + grey dek,
% all left-aligned to a shared margin.  \fighead{x}{y (top-left)}{KICKER}{Declarative title}{dek}
\newcommand{\fighead}[5]{%
  \draw[claret, line width=2pt] (#1,#2) -- ++(0.62,0);
  \node[anchor=north west, inner sep=0pt, font=\scriptsize\bfseries, text=claret]
       (fhk) at ($(#1,#2)+(0,-0.18)$) {\textls[140]{#3}};
  \node[anchor=north west, inner sep=0pt, font=\rmfamily\LARGE\bfseries, text=ink]
       (fht) at ($(fhk.south west)+(0,-0.13)$) {#4};
  \node[anchor=north west, inner sep=0pt, font=\small, text=sub]
       (fhd) at ($(fht.south west)+(0,-0.13)$) {#5};
}
% bottom-left source / credit line (the FT/Nature/Quanta honesty finish). \figsource{x}{y}{text}
\newcommand{\figsource}[3]{%
  \node[anchor=north west, inner sep=0pt, font=\scriptsize, text=mute, align=left] at (#1,#2) {#3};
}
% a filled tag set as a card's first line. \tagline{fill colour}{TEXT}  (shared across figures)
\newcommand{\tagline}[2]{{\colorbox{#1}{\textcolor{white}{\scriptsize\bfseries #2}}}\\[4pt]}

% ---- plain-language causal-chain flow (general audience, icon-led) -----------
\tikzset{
  stagecircle/.style={circle, minimum size=1.55cm, inner sep=0pt, text=white,
                      font=\LARGE, line width=0pt},
  stagelabel/.style={font=\small, text=ink, align=center, text width=3.1cm},
  flowedge/.style={-{Stealth[length=3.4mm,width=3mm]}, line width=1.7pt, sub},
  flowtitle/.style={font=\Large\bfseries, text=ink},
  flowcap/.style={font=\small, text=sub, align=center},
}
% \stage{id}{x}{y}{fillcolor}{icon}{label}
\newcommand{\stage}[6]{%
  \node[stagecircle, fill=#4] (#1) at (#2,#3) {#5};
  \node[stagelabel, below=2.5mm of #1] {#6};
}

% a small "data" sparkline card that shows the by-X profile a channel uses, with a
% source tag -- the visible link between a mechanism and its model effect.
% \datacard{x}{y}{title}{x-axis label}{source tag}{colour}{plot coords}
% coords are relative to (x,y-0.05), spanning x in [-1.1,1.1], y in [-0.4,0.45].
\newcommand{\datacard}[7]{%
  \node[draw=grid, rounded corners=3pt, fill=white, minimum width=2.8cm,
        minimum height=1.8cm, inner sep=0pt] at (#1,#2) {};
  \node[font=\tiny\bfseries, text=sub, anchor=north] at (#1,#2+0.80) {#3};
  \begin{scope}[shift={(#1,#2-0.05)}]
    \draw[grid, line width=0.6pt] (-1.12,-0.42) -- (1.12,-0.42);
    \draw[#6, line width=1.4pt, line cap=round] plot[smooth] coordinates {#7};
    \node[font=\tiny, text=mute, anchor=north west] at (-1.08,-0.42) {#4};
  \end{scope}
  \node[font=\tiny\bfseries, text=white, fill=#6, rounded corners=1.5pt, inner xsep=2.5pt,
        inner ysep=0.8pt, anchor=south east] at (#1+1.3,#2-0.86) {#5};
}


\tikzset{
  pivotc/.style={card, draw=clews, line width=1.6pt, fill=clews!6, text width=3.5cm, font=\footnotesize},
  arm/.style={card, draw=grid, line width=0.9pt, fill=white, text width=4.7cm, font=\footnotesize},
  hhc/.style={card, draw=ink, line width=1.5pt, fill=ink!3, text width=2.7cm, align=center, font=\footnotesize},
  outc/.style={card, draw=grid, fill=white, text width=3.0cm, font=\scriptsize, align=left},
  edgefaith/.style={flow, draw=ink!78},
  edgeill/.style={flow, draw=ink!78, dash pattern=on 1pt off 2.4pt},
  edgeloss/.style={flow, draw=loss},
  modarrow/.style={-{Stealth[length=2mm,width=1.7mm]}, draw=mute, line width=0.8pt,
                   dash pattern=on 1.6pt off 1.6pt},
  modchip/.style={draw=mute, dash pattern=on 1.6pt off 1.6pt, rounded corners=3pt, fill=white,
                  font=\scriptsize, text=sub, align=center, inner sep=4pt, text width=3.2cm},
}
\newcommand{\modelbox}[1]{\\[4pt]{\scriptsize\textcolor{ogc}{\textbf{In the model:}}\ #1}}

\begin{document}
\begin{tikzpicture}[every node/.style={outer sep=0pt}]

\fighead{-1.7}{8.3}{THE ENERGY-PRICE FORK}{One electricity price hits the economy through two channels at once}{producers' costs rise, and households pay more directly --- both reach households}

% ---------- nodes ----------
\node[pivotc] (piv) at (0,0) {%
  {\bfseries\color{clews}Electricity costs more to supply}\\[2pt]
  {\scriptsize\color{sub}the price the CLEWS energy model returns --- the cost of supplying one more unit of electricity}};

\node[arm] (prod) at (6.3,3.1) {%
  \tagline{ink}{CALIBRATED FROM DATA}%
  {\bfseries Producers' costs rise}\ {\bfseries\color{loss}($-$)}\\[2pt]
  Every industry that uses electricity faces higher costs, so it produces less from the same inputs and its goods cost more.%
  \modelbox{the economy model tracks each industry's labour and capital, not the energy it buys --- so this is approximated: each industry's productivity is trimmed by how much electricity it uses.}};

\node[arm] (dem) at (6.3,-3.1) {%
  \tagline{ink}{REPRESENTED DIRECTLY}%
  {\bfseries Households pay more directly}\\[2pt]
  Households pay more for the energy they buy, so they use less and spend more on a necessity.%
  \modelbox{this is how the economy model naturally represents an energy price. When energy is set as a necessity and the revenue is recycled, the burden is regressive --- shown by a dedicated incidence run.}};

\node[hhc] (hh) at (12.0,0) {%
  {\bfseries Households}\\[2pt]{\scriptsize\color{sub}both channels reach households}};

\node[outc] (o1) at (16.0,2.6) {{\bfseries\color{ogc}GDP falls}\ {\color{loss}$\downarrow$}\\{\scriptsize\color{sub}mainly the producers'-cost channel}};
\node[outc] (o2) at (16.0,0)   {{\bfseries\color{ogc}Real incomes fall}\ {\color{loss}$\downarrow$}};
\node[outc] (o3) at (16.0,-2.6){{\bfseries\color{ogc}Who bears it}\\{\scriptsize\color{sub}regressive when energy is treated as a necessity}};

% ---------- edges (auto-attach to borders) ----------
\draw[edgeill]   (piv.east) to[out=58,in=205] node[pos=0.5,coordinate](pe){} (prod.west);
\draw[edgefaith] (piv.east) to[out=-58,in=155] (dem.west);
\draw[edgeill]   (prod.east) to[out=-15,in=125] (hh.north west);
\draw[edgefaith] (dem.east) to[out=15,in=-125] (hh.south west);
\draw[edgeloss]  (hh.east) to[out=32,in=180] (o1.west);
\draw[edgeloss]  (hh.east) to[out=0,in=180]  (o2.west);
\draw[edgefaith] (hh.east) to[out=-32,in=180] (o3.west);

% ---------- sizing factor landing ON the pivot->prod edge ----------
\node[modchip] (mod) at (0.8,4.3) {how much electricity each industry uses};
\draw[modarrow] (mod.east) to[out=-30,in=150] (pe);

% ---------- legend ----------
\begin{scope}[shift={(-0.1,-5.7)}]
  \node[font=\scriptsize\bfseries, text=ink, anchor=west] at (0,0) {How each channel is represented};
  \draw[edgefaith] (5.0,0) -- (5.9,0);
  \node[font=\scriptsize, text=sub, anchor=west] at (5.95,0) {represented directly --- the model's own mechanism (size reliable)};
  \draw[edgeill] (0,-0.7) -- (0.9,-0.7);
  \node[font=\scriptsize, text=sub, anchor=west] at (0.95,-0.7) {calibrated from data --- a stand-in for the energy in production (size approximate)};
  \node[text=loss, font=\bfseries] at (0.2,-1.4) {$\downarrow$};
  \node[font=\scriptsize, text=sub, anchor=west] at (0.55,-1.4) {falls / contractionary};
  \node[modchip, scale=0.82, text width=2.0cm] at (5.9,-1.4) {sizing factor};
  \node[font=\scriptsize, text=sub, anchor=west] at (7.0,-1.4) {scales the channel it points at};
\end{scope}

\node[font=\scriptsize, text=mute, align=left, anchor=north west, text width=16.8cm] at (-0.1,-7.6)
  {Colour shows which model each element comes from (teal = the CLEWS energy model, blue = the OG-Core economy),
   not its sign --- direction is shown by the arrows and line style, so the figure reads the same in greyscale.
   The linked run combines these two channels on disjoint flows --- electricity's own use is carried by the
   household wedge, so the price is not counted twice.};

\end{tikzpicture}
\end{document}
